\documentclass[11pt]{amsart}
\usepackage[T1]{fontenc}
\usepackage[utf8]{inputenc}
\usepackage{amssymb}
\usepackage[margin=1.05in]{geometry}
\usepackage[hidelinks]{hyperref}
\usepackage{booktabs}
\usepackage{array}
\usepackage{tikz-cd}
\usepackage{datetime}
\numberwithin{equation}{section}
\newtheorem{theorem}{Theorem}[section]
\newtheorem{proposition}[theorem]{Proposition}
\newtheorem{lemma}[theorem]{Lemma}
\newtheorem{corollary}[theorem]{Corollary}
\theoremstyle{definition}
\newtheorem{definition}[theorem]{Definition}
\newtheorem{axiom}[theorem]{Axiom}
\theoremstyle{remark}
\newtheorem{remark}[theorem]{Remark}
\newcommand{\code}[1]{\texttt{#1}}
\newcommand{\Cat}{\mathsf{Cat}}
\newcommand{\Map}{\mathrm{Map}}
\newcommand{\Fun}{\mathrm{Fun}}
\newcommand{\Ob}{\mathrm{Ob}}
\newcommand{\Mor}{\mathrm{Mor}}
\newcommand{\Hom}{\mathrm{Hom}}
\newcommand{\Iso}{\mathrm{Iso}}
\newcommand{\Ch}{\mathrm{Ch}}
\newcommand{\dom}{\mathrm{dom}}
\newcommand{\cod}{\mathrm{cod}}
\newcommand{\id}{\mathrm{id}}
\newcommand{\II}{\mathbb{I}}
\newcommand{\one}{\mathbf{1}}
\newcommand{\isContr}{\mathrm{isContr}}
\newcommand{\isEquiv}{\mathrm{isEquiv}}
\AddToHook{shipout/foreground}{%
 \put(0.5\paperwidth,-\paperheight+0.55cm){\makebox[0pt]{\tiny\ttfamily
 unit/doc/\jobname.tex\ ---\ compiled \today\ \currenttime}}}
\title[The Segal condition and fully faithful functors]{The Segal condition
for all $n$, composition, and fully faithful functors\\ from the
Sattler--W\"arn axioms, and their formalisation in Unit}
\date{31 October 2026}
\begin{document}
\sloppy
\begin{abstract}
We work in univalent type theory with an axiomatised wild category
$\Cat$ of higher categories, following W\"arn's notes and the axioms of
Sattler and W\"arn. From the Segal axiom and the fact that $\Delta^{n+1}$
is obtained from $\Delta^n$ by gluing an edge (Lemma 26 of W\"arn's
notes) we prove the Segal condition for every $n$: a functor
$\Delta^n\to C$ is the same as a chain of $n$ composable morphisms of
$C$. Composition of morphisms is defined as the centre of a contractible
type of fillers, and the unit and associativity laws follow from the
cases $n=2$ and $n=3$. We then prove W\"arn's results 33--39 on fully
faithful functors: they reflect invertibility, are embeddings on
objects, and are isomorphisms when they are surjective on objects. All
of this uses only two coherences of $\Cat$ beyond its associativity and
unit laws: the pentagon and the triangle. The last section explains
how the proofs are formalised in the proof checker Unit.
\end{abstract}
\maketitle

\section{Introduction}

W\"arn's notes \emph{Higher-categorical headcanon} and the axioms of
Sattler and W\"arn present $(\infty,1)$-category theory on top of
homotopy type theory, after Cisinski, Cnossen, Nguyen and Walde. One
postulates a wild category $\Cat$, whose objects are called categories
and whose morphisms are called functors, with a few universal
properties: a terminal category, pullbacks, exponentials, the walking
arrow $\II$, and the Segal axiom. Everything internal to a category
$C$ (its objects, morphisms, composition and the laws of composition)
is then constructed. Being a wild category is a matter of coherences:
the associativity of functor composition holds up to a path, these
paths satisfy the pentagon up to a 2-path, and so on. W\"arn treats the
whole tower as an axiom schema and expects that every concrete argument
uses only finitely many levels.

This note proves the following.
\begin{enumerate}
\item The Segal condition for every $n$ (Theorem~\ref{thm:segal}). The
  case $n=2$ is an axiom; the other cases follow from it and from
  W\"arn's Lemma 26.
\item Composition of morphisms with its unit and associativity laws,
  and the fact that functors preserve composition and identities
  (Section~\ref{sec:comp}).
\item W\"arn's results 33--39 on invertible morphisms and fully
  faithful functors (Section~\ref{sec:ff}), with the converse of
  Remark~39.
\end{enumerate}
We also record which coherences each result uses
(Section~\ref{sec:levels}). Composition itself uses only the
associativity and unit paths of $\Cat$. The pentagon and the triangle
enter because a functor acts on a morphism together with the
identifications of its endpoints, and the objects of a category do not
form a set. No higher coherence is used.

\section{The axioms}\label{sec:axioms}

We work in dependent type theory with univalence; function
extensionality follows. Equality means the identity type, and
$p\cdot q$ is the composite of paths.

\subsection{Level 1: the wild category}
We are given a type $\Cat$ of \emph{categories}, for $A,B:\Cat$ a type
$\Map(A,B)$ of \emph{functors}, composition $g\circ f$, identities
$1_A$, and paths
\[
\alpha_{h,g,f}:(h\circ g)\circ f=h\circ(g\circ f),\qquad
\lambda_f:1\circ f=f,\qquad \rho_f:f\circ 1=f.
\]
The types $\Map(A,B)$ are not assumed to be sets.
A functor $f:A\to B$ is an \emph{isomorphism} if
$f\circ-:\Map(X,A)\to\Map(X,B)$ is an equivalence for every $X$. This is
a proposition, and an isomorphism has a two-sided inverse.

\subsection{Level 2: pentagon and triangle}
We are given the pentagon, a path between the two ways of reassociating
$((f\circ g)\circ h)\circ e$ into $f\circ(g\circ(h\circ e))$, and the
triangle relating $\alpha_{g,1,f}$, $\rho_g$ and $\lambda_f$. The other
two triangles and Mac Lane's coherence theorem (any two parallel
composites of $\alpha,\lambda,\rho$ and their inverses are equal)
follow, as usual.

\subsection{Structure}
The following axioms only mention level-1 data.
\begin{itemize}
\item $\Cat$ is univalent (Axiom 3), and has a terminal category $\one$
  (Axiom 4), pullbacks (Axiom 13) and exponentials $\Fun(A,B)$ with
  $\Map(X,\Fun(A,B))\simeq\Map(X\times A,B)$ (Axiom 15).
\item The objects of $C$ are $\Ob C=\Map(\one,C)$, and a functor acts on
  them by composition.
\item There is a category $\II$ with $\Ob\II\simeq\{0,1\}$ (Axiom 8).
  Writing $\Mor C=\Map(\II,C)$, $\dom m=m\circ 0$ and $\cod m=m\circ 1$,
  every morphism of $\II$ goes from its domain to an object that is at
  least as large. Moreover $\Map(C,\II)$ is equivalent to the type of
  maps $\Ob C\to\{0,1\}$ that are monotone along every morphism
  (Axiom 18).
\item The Segal axiom (Axiom 21): with $\Delta^0=\one$ and
  $\Delta^{n+1}=\Fun(\Delta^n,\II)$ (Definition 19), the map
  \[
  \Map(\Delta^2,C)\to(f,g:\Mor C)\times(\cod f=\dom g),\qquad
  s\mapsto(s\circ d_{01},\,s\circ d_{12})
  \]
  is an equivalence. Here $d_{01},d_{12}$ are the two spine edges of
  $\Delta^2$.
\end{itemize}
The \emph{mapping space} is
\[
\Hom_C(x,y)=(m:\Mor C)\times(\dom m=x)\times(\cod m=y).
\]
Since $\Ob C$ is not a set, a path in $\Hom_C(x,y)$ is a path $m=m'$ of
morphisms together with two 2-cells relating the endpoint proofs. The
identity of $c$ is $\id_c=c\circ{!}$ with ${!}:\II\to\one$, and its
endpoint proofs use that $\Map(\one,\one)$ is contractible. A functor
$F:A\to B$ acts on $\Hom$:
\[
F(m,p,q)=(F\circ m,\ \alpha_{F,m,0}\cdot \mathrm{ap}_F\,p,\
\alpha_{F,m,1}\cdot\mathrm{ap}_F\,q):\Hom_B(Fx,Fy).
\]
Two more axioms concern invertible morphisms; they are stated in
Section~\ref{sec:ff}, once composition has been defined.

\subsection{Simplices}
From Axiom 18 one gets $\Ob\Delta^n\simeq[n]=\{0<1<\dots<n\}$, and
functors into $\Delta^n$ are monotone maps on objects (Construction 20).
Every $\Delta^n$ and every $\Fun(X,\II)$ is \emph{posetal}: its mapping
spaces are propositions, its objects form a set, and a functor into it
is determined by its action on objects. Consequently $\Map(A,\Delta^n)$
is a set. We write $\mathrm{top}_n$ for the last vertex of $\Delta^n$,
$\phi_n:\Delta^n\to\Delta^{n+1}$ for the face that omits the last
vertex, and $\ell_n:\Hom_{\Delta^{n+1}}(\phi_n(\mathrm{top}_n),
\mathrm{top}_{n+1})$ for the last edge.

\section{Gluing an edge to a simplex}\label{sec:lemma26}

The \emph{gluing square} $S_n$ is
\[
\begin{tikzcd}
\one \arrow[r,"0"] \arrow[d,"\mathrm{top}_n"'] & \Delta^1 \arrow[d,"\ell_n"]\\
\Delta^n \arrow[r,"\phi_n"'] & \Delta^{n+1}
\end{tikzcd}
\]
where the right edge is $\ell_n$ precomposed with the isomorphism
$\Delta^1\simeq\II$. A commuting square is a \emph{pushout} when, for
every $X$, maps out of its corner $\Delta^{n+1}$ correspond, by
precomposition, to cones over its span.

\begin{theorem}[W\"arn, Lemma 26]\label{thm:lemma26}
For every $n$ the square $S_n$ is a pushout.
\end{theorem}

\begin{proof}
For $n=0$ the left edge is an identity. For $n=1$ the cones over the
span are composable pairs, through $\Delta^1\simeq\II$, and the
statement is the Segal axiom. For the induction step we use three
facts.
\begin{enumerate}
\item $-\times\II$ preserves pushouts, since it has the right adjoint
  $\Fun(\II,-)$. So $S_{n+1}\times\II$ is a pushout.
\item $S_{n+2}$ is a retract of $S_{n+1}\times\II$. At the corners
  $\Delta^{k+1}$ the section is $\psi\mapsto(\psi\circ\phi_k,\
  \psi(\mathrm{top}_{k+1}))$, which restricts a functor
  $\psi:\Delta^{k+1}\to\II$ to the face and evaluates it at the top. The
  retraction is induced by a monotone retraction of posets
  $[k+1]\times\{0,1\}\to[k+2]$. At the corners $\one$ and $\Delta^1$
  the section is $x\mapsto(x,1)$ and the retraction is the first
  projection. That the sides commute is checked on objects, as all the
  corners are posetal.
\item A retract of a pushout square is a pushout square, provided the
  types of functors between corners are sets. Here they are, since the
  corners are posetal. Under that hypothesis a morphism of squares
  needs no coherence between its faces, and the comparison map of the
  retract is a retract of an equivalence. \qedhere
\end{enumerate}
\end{proof}

The same holds with $\II$ in place of $\Delta^1$. That square is a
retract of $S_n$ through the isomorphism $\Delta^1\simeq\II$. So for
every $n$ and $C$ the map
\begin{equation}\label{eq:pushout}
\Map(\Delta^{n+1},C)\to(a:\Map(\Delta^n,C))\times(m:\Mor C)\times
\bigl(a(\mathrm{top}_n)=\dom m\bigr),
\end{equation}
which sends $b$ to $b\circ\phi_n$ and $b\circ\ell_n$ together with the
path given by the commuting square, is an equivalence.

\section{The Segal condition for all $n$}\label{sec:segal}

\begin{definition}
The type $\Ch_n(x)$ of \emph{chains} of $n$ composable morphisms ending
at $x:\Ob C$ is defined by recursion on $n$:
\[
\Ch_0(x)=\one,\qquad
\Ch_{n+1}(x)=(y:\Ob C)\times\Ch_n(y)\times\Hom_C(y,x).
\]
So a chain of length $n+1$ is a chain of length $n$ with a morphism
glued at its last vertex. A functor $F$ maps chains to chains, vertex
by vertex and edge by edge. The \emph{universal chain}
$u_n:\Ch_n(\mathrm{top}_n)$ of $\Delta^n$ is defined by $u_0=\star$
and $u_{n+1}=(\phi_n(\mathrm{top}_n),\ \phi_n(u_n),\ \ell_n)$. The
\emph{spine} of $s:\Delta^n\to C$ is $s(u_n):\Ch_n(s(\mathrm{top}_n))$.
\end{definition}

\begin{lemma}\label{lem:homfunct}
For $G:A\to B$, $F:B\to C$ and $h:\Hom_A(x,y)$, the morphism
$(F\circ G)(h)$, with its endpoints moved along $\alpha$, is equal to
$F(G(h))$. Consequently $(F\circ G)(c)=F(G(c))$ for chains $c$, over
$\alpha$ at the last vertex. Also $F(\id_x)=\id_{Fx}$, and
$m:\Mor C$ is the image of the generating arrow of $\II$.
\end{lemma}

\begin{proof}
By path induction on the endpoint proofs of $h$, we may take
$h=(m,\mathrm{refl},\mathrm{refl})$. The morphisms agree by
$\alpha_{F,G,m}$. The two endpoint 2-cells are the pentagon at
$(F,G,m,0)$ and at $(F,G,m,1)$. For chains, induct on the length. For
the identity, the morphisms agree by $\alpha_{F,c,!}^{-1}$, and the
endpoint 2-cells agree by the pentagon at $(F,c,!,0)$, naturality of
$\alpha$, and the triangle. The last statement is the triangle at
$(0,m)$ and $(1,m)$.
\end{proof}

\begin{theorem}[Segal condition]\label{thm:segal}
For every $n$ and every category $C$, the spine map
\[
\Map(\Delta^n,C)\longrightarrow (x:\Ob C)\times\Ch_n(x),\qquad
s\mapsto (s(\mathrm{top}_n),\,s(u_n))
\]
is an equivalence.
\end{theorem}

\begin{proof}
By induction on $n$. For $n=0$, both sides are $\Ob C$. For $n+1$, the
spine map is homotopic to the composite of three equivalences:
\[
\Map(\Delta^{n+1},C)\ \simeq\
(a,m,p)\ \simeq\ ((x,c),m,\,x=\dom m)\ \simeq\
(x:\Ob C)\times\Ch_{n+1}(x).
\]
The first is \eqref{eq:pushout}. The second applies the induction
hypothesis to $a$. The third glues: $((x,c),m,p)\mapsto(\cod m,\,x,\,c,\,
(m,p^{-1},\mathrm{refl}))$, an equivalence by contracting the
singleton $(\cod m,\mathrm{refl})$. To see that the composite is the
spine map, take $b:\Delta^{n+1}\to C$.
\begin{itemize}
\item The last vertices agree, since
  $\cod(b\circ\ell_n)=b(\mathrm{top}_{n+1})$.
\item The chains agree: the spine of $b\circ\phi_n$ is
  $(b\circ\phi_n)(u_n)=b(\phi_n(u_n))$ by Lemma~\ref{lem:homfunct}.
\item The last morphisms agree by a reassociation of the commuting
  square.\qedhere
\end{itemize}
\end{proof}

\begin{corollary}
The type of $n$-simplices of $C$ with a given spine is contractible.
\end{corollary}

\section{Composition and its laws}\label{sec:comp}

Let $\mathrm{Pair}_C$ be the type of composable pairs with their
vertices, $(z,y,x)$ with $f:\Hom(z,y)$ and $g:\Hom(y,x)$. The
\emph{boundary} of $s:\Delta^2\to C$ is the image under $s$ of the three
edges $d_{01},d_{12},d_{02}$ of $\Delta^2$ with their endpoints. Its
first two edges form a pair $\partial s:\mathrm{Pair}_C$.

\begin{lemma}\label{lem:seg2}
$\partial:\Map(\Delta^2,C)\to\mathrm{Pair}_C$ is an equivalence. Its
fibres, the \emph{fillers} of a pair, are contractible.
\end{lemma}

\begin{proof}
$\partial$ is the Segal comparison followed by the equivalence
$(f,g,p)\mapsto(\dom f,\cod f,\cod g,(f,\mathrm{refl},\mathrm{refl}),
(g,p^{-1},\mathrm{refl}))$. The inverse of that equivalence forgets the
vertices, which works by contracting the endpoint proofs. The two
constructions agree on the path $\cod(s\,d_{01})=\dom(s\,d_{12})$ by
the groupoid laws.
\end{proof}

Lemma~\ref{lem:seg2} uses only the Segal axiom, not
Theorem~\ref{thm:segal}. The spine of a 2-simplex, read as a pair, is
exactly $\partial s$. This is because the spine of $\Delta^2$ consists
of the edges $d_{01}$ and $d_{12}$, and pairs in the posetal $\Delta^2$
are determined by their vertices. So Lemma~\ref{lem:seg2} is the case
$n=2$ of Theorem~\ref{thm:segal}.

\begin{definition}[Construction 22]
For a pair $(f,g)$, the composite $g\circ f:\Hom(z,x)$ is the third
edge of the centre of its type of fillers, transported along the
identification of the centre's first two edges with $(f,g)$.
\end{definition}

Composition is characterised by two properties.
\begin{enumerate}
\item If $s:\Delta^2\to C$ has boundary $(f,g,h)$, then $h=g\circ f$.
  This holds because $(s,\ldots)$ is a filler and fillers form a
  contractible type.
\item Every pair $(f,g)$ has a 2-simplex with boundary
  $(f,g,g\circ f)$.
\end{enumerate}
In particular, on the boundary of any 2-simplex, $g\circ f$ has the
same underlying morphism as the long edge of the witness given by the
Segal axiom.

\begin{proposition}\label{prop:funct}
For a functor $F:A\to C$: $F(g\circ f)=Fg\circ Ff$ and $F(\id_x)=\id_{Fx}$.
\end{proposition}

\begin{proof}
Let $s$ have boundary $(f,g,g\circ f)$. By Lemma~\ref{lem:homfunct}
applied to each edge, $F\circ s$ has boundary $(Ff,Fg,F(g\circ f))$.
Property (1) of composition gives the first statement. The second is
in Lemma~\ref{lem:homfunct}.
\end{proof}

\begin{proposition}[Construction 23]
$h\circ\id_a=h$ and $\id_b\circ h=h$ for $h:\Hom_C(a,b)$.
\end{proposition}

\begin{proof}
By path induction on the endpoint proofs, take $h=(m,\mathrm{refl},
\mathrm{refl})$ with $m:\II\to C$. In $\II$ the pair $(\id_0,\gamma)$,
where $\gamma$ is the generating arrow, has a filler $\sigma$ by
Lemma~\ref{lem:seg2}. Its long edge is $\gamma$, since $\Hom_\II(0,1)$
is a proposition. By Lemma~\ref{lem:homfunct}, $m\circ\sigma$ has
boundary $(m\,\id_0,\,m\gamma,\,m\gamma)=(\id_a,h,h)$. So $h\circ\id_a=h$
by property (1). The other law uses the pair $(\gamma,\id_1)$.
\end{proof}

\begin{proposition}[Construction 27]
$(k\circ g)\circ f=k\circ(g\circ f)$ for composable $f,g,k$.
\end{proposition}

\begin{proof}
By Theorem~\ref{thm:segal} at $n=3$, the type of triples of composable
morphisms with their vertices is equivalent to $\Map(\Delta^3,C)$,
through the spine. It therefore suffices to prove the law for the
spine of some $T:\Delta^3\to C$, that is, for $(Tk_{01},Tk_{12},Tk_{23})$
where $k_{i,i+1}$ is the spine of $\Delta^3$. By
Proposition~\ref{prop:funct} both sides are images under $T$:
\[
(Tk_{23}\circ Tk_{12})\circ Tk_{01}=T((k_{23}\circ k_{12})\circ k_{01}),\qquad
Tk_{23}\circ(Tk_{12}\circ Tk_{01})=T(k_{23}\circ(k_{12}\circ k_{01})).
\]
The two composites in $\Delta^3$ are equal, because $\Delta^3$ is
posetal and its mapping spaces are propositions.
\end{proof}

\section{Invertible morphisms and fully faithful functors}\label{sec:ff}

\begin{definition}[Definitions 33 and 34]
A morphism $f:\Hom_C(a,b)$ is \emph{invertible} if, for every $x$,
postcomposition $f\circ-:\Hom(x,a)\to\Hom(x,b)$ is an equivalence. This
is a proposition. Let $\Iso_C(a,b)$ be the type of invertible
morphisms. A functor $F$ is \emph{fully faithful} if each action
$\Hom_A(a_0,a_1)\to\Hom_B(Fa_0,Fa_1)$ is an equivalence.
\end{definition}

\begin{proposition}\label{prop:inv}
$f$ is invertible if and only if it has a two-sided inverse.
Identities are invertible, and functors preserve invertibility.
\end{proposition}

\begin{proof}
If $g$ is a two-sided inverse, then $g\circ-$ is inverse to $f\circ-$,
by associativity and the unit laws. Conversely, let $g$ be the preimage
of $\id_b$ under $f\circ-$. Then $f\circ g=\id_b$. Moreover
$f\circ(g\circ f)=(f\circ g)\circ f=f=f\circ\id_a$, and $f\circ-$ is
injective, so $g\circ f=\id_a$. Identities have themselves as inverses.
Functors preserve two-sided inverses by Proposition~\ref{prop:funct}.
\end{proof}

\begin{lemma}[Lemma 35]
A fully faithful functor reflects invertibility.
\end{lemma}

\begin{proof}
By Proposition~\ref{prop:funct}, $F\circ(f\circ-)=(Ff\circ-)\circ F$ as
maps $\Hom(x,a)\to\Hom(Fx,Fb)$. If $Ff$ is invertible and $F$ is fully
faithful, the right side is an equivalence. So $f\circ-$ is one by
cancellation.
\end{proof}

\begin{axiom}[Rezk, Axiom 36]\label{ax:rezk}
For every $a:\Ob C$, the type $(b:\Ob C)\times\Iso_C(a,b)$ is
contractible.
\end{axiom}

\begin{remark}\label{rem:rezk}
W\"arn states Axiom 36 as follows: $c\mapsto\id_c$, $\Ob C\to\Mor C$, is
an embedding whose image contains the invertible morphisms. The two
forms are equivalent. Since $\Mor C\simeq(a,b:\Ob C)\times\Hom(a,b)$ by
contracting the endpoint proofs, and invertibility is a proposition,
W\"arn's form says that $\iota:c\mapsto(c,c,\id_c)$ is an equivalence
$\Ob C\simeq(a:\Ob C)\times(b:\Ob C)\times\Iso(a,b)$. This holds
because a surjective embedding is an equivalence, and identities are
invertible. As $\iota$ lies over the identity of $\Ob C$, it is an
equivalence exactly when each map $\one\to(b:\Ob C)\times\Iso(a,b)$ is
one, which is Axiom~\ref{ax:rezk}.
\end{remark}

By the fundamental theorem of identity types, Axiom~\ref{ax:rezk} is
equivalent to the following: the map $\mathrm{idtoiso}:(a=b)\to
\Iso(a,b)$, which transports $\id_a$, is an equivalence.

\begin{lemma}[Lemma 37]
If $F$ is fully faithful, then $\Ob F:\Ob A\to\Ob B$ is an embedding.
\end{lemma}

\begin{proof}
The action of $F$ on invertible morphisms, $\Iso(a_0,a_1)\to
\Iso(Fa_0,Fa_1)$, is an equivalence. It is the action on $\Hom$, an
equivalence, over a map between propositions that goes both ways by
Proposition~\ref{prop:inv} and Lemma 35. By path induction and
$F(\id)=\id$, we have $\mathrm{idtoiso}(\mathrm{ap}_F\,p)=
F(\mathrm{idtoiso}\,p)$. So
$\mathrm{ap}_F=\mathrm{idtoiso}^{-1}\circ F\circ\mathrm{idtoiso}$ is an
equivalence.
\end{proof}

\begin{axiom}[$\II$ detects equivalences, Axiom 38]
If $F\circ-:\Map(\II,A)\to\Map(\II,B)$ is an equivalence, then $F$ is an
isomorphism.
\end{axiom}

\begin{theorem}[Remark 39 and its converse]
A functor $F:A\to B$ is an isomorphism if and only if it is fully
faithful and surjective on objects. Here surjective means that every
fibre of $\Ob F$ is merely inhabited.
\end{theorem}

\begin{proof}
Suppose $F$ is fully faithful and surjective. By Lemma 37, $\Ob F$ is
an embedding. Being surjective, it is an equivalence. Through
$\Mor C\simeq(a,b)\times\Hom(a,b)$, postcomposition $F\circ-$ on
morphisms becomes $((a,b),h)\mapsto((Fa,Fb),Fh)$. This is a fibrewise
equivalence followed by a change of base along the equivalence
$\Ob F\times\Ob F$, so it is an equivalence. Axiom 38 then shows that
$F$ is an isomorphism.

Conversely, suppose $F$ is an isomorphism. Postcomposition at $X=\one$
shows that $\Ob F$ is an equivalence, hence surjective. Postcomposition
at $X=\II$ shows that the total map above is an equivalence. After
cancelling the change of base, the map over $\Ob A\times\Ob A$ with
fibres $h\mapsto Fh$ is an equivalence of total spaces. So it is a
fibrewise equivalence, which says that $F$ is fully faithful.
\end{proof}

\section{Which coherences are used}\label{sec:levels}

\begin{center}
\begin{tabular}{>{\raggedright}p{6.6cm}l>{\raggedright\arraybackslash}p{4.6cm}}
\toprule
Result & Level & Coherence used\\
\midrule
objects, morphisms, $\Delta^n$, Construction 20, posetality & 1 & --\\
$\II$-faces of $\Delta^2$, retract $S_{n+2}<S_{n+1}\times\II$ & 1 & --\\
composition, its characterisation (Lemma~\ref{lem:seg2}) & 1 & --\\
definitions of invertible, fully faithful; Axiom 36 & 1 & --\\
Lemma~\ref{lem:homfunct} & 2 & pentagon, triangle\\
Theorem~\ref{thm:lemma26} (Lemma 26) & 2 & pentagon (cones, products)\\
Theorem~\ref{thm:segal} (Segal for all $n$) & 2 & Lemma 26, Lemma~\ref{lem:homfunct}\\
unit laws, associativity, Proposition~\ref{prop:funct} & 2 & Lemma~\ref{lem:homfunct}\\
results 35--39 & 2 & through the above\\
\bottomrule
\end{tabular}
\end{center}

\noindent
Nothing uses a coherence of level 3. Level 2 enters in two places:
\begin{itemize}
\item in the algebra of cones and products used for Lemma 26;
\item in the action of functors on morphisms together with their
  endpoint identifications (Lemma~\ref{lem:homfunct}).
\end{itemize}
The second is needed because $\Ob C$ is not a set.

\section{The formalisation}\label{sec:formal}

Everything above is checked in Unit, a proof checker for dependent type
theory (\url{github.com/coquand/unit}), in the directory
\code{unit/sag/hcat}. Sections~\ref{sec:axioms}--\ref{sec:lemma26}
port an Agda proof of Lemma 26 (\code{github.com/coquand/basic-higher-category}).
Sections~\ref{sec:segal}--\ref{sec:levels} are new.

\subsection{Axioms as nested segments}
Unit has \emph{segments}: modules whose parameters are assumptions,
and which can be nested. The root file \code{All.seg} is laid out as
\begin{center}
\code{Theory} $\supset$ \code{Wild} $\supset$ \code{Structure}
$\supset$ \code{Coh2}.
\end{center}
\begin{itemize}
\item \code{Theory} takes the parameters \code{univ} (univalence),
  \code{images} and \code{surjections} (replacement). Function
  extensionality is derived.
\item \code{Wild} takes the level-1 data of $\Cat$.
\item \code{Structure} takes the axioms of Section~\ref{sec:axioms}
  together with \code{rezk} and \code{I-detects-equiv}.
\item \code{Coh2} takes \code{pentagonator} and \code{triangle2}.
\end{itemize}
Each file is imported in the outermost segment whose parameters it
uses, so the level column of Section~\ref{sec:levels} can be read from
the layout. The definition of $\Cat$ has no judgemental shortcuts: the
laws of functor composition are paths, and the two triangles that are
not parameters, as well as Mac Lane coherence, are proved from the
pentagon and \code{triangle2}. The
formalisation uses no new axioms beyond the \code{Theory} parameters
\code{univ}, \code{images}, \code{surjections} and the parameters of
\code{Wild}, \code{Structure} and \code{Coh2}.

\subsection{Where the results are}
\begin{center}
\begin{tabular}{>{\raggedright}p{5.2cm}>{\raggedright\arraybackslash}p{8.6cm}}
\toprule
Result & File and name\\
\midrule
Theorem~\ref{thm:lemma26} & \code{Lemma26.lemma26-square}\\
the square with $\II$, \eqref{eq:pushout} & \code{SpineSquare.pushout-is-equiv}\\
chains, $u_n$, spine & \code{Chains.HChain}, \code{univ-chain}, \code{spine-total}\\
Lemma~\ref{lem:homfunct} & \code{HomFunctor.Hom-map-of-comp}, \code{Hom-map-id},
  \code{Hom-map-gen}; \code{ChainFunctor.Chain-map-of-comp}\\
Theorem~\ref{thm:segal}, corollary & \code{SegalAll.segal-all}, \code{filler-contr}\\
Lemma~\ref{lem:seg2}, composition & \code{Composition.segal2-equiv}, \code{composite} (locked),
  \code{composite-witness}, \code{composite-filler}, \code{hcomp}\\
spine at $n=2$ is $\partial$ & \code{SpineTwo.spine-two}\\
agreement with the Segal witness & \code{Composition.composite-agrees}\\
Proposition~\ref{prop:funct}, unit laws & \code{CompLaws.Hom-map-hcomp}, \code{hcomp-unit-r}, \code{hcomp-unit-l}\\
associativity & \code{Associativity.hcomp-assoc}\\
invertible, fully faithful & \code{Invertible.is-inv}, \code{is-inv-prop}, \code{ObIso}, \code{is-ff}\\
Proposition~\ref{prop:inv} & \code{FullyFaithful.has-inv-is-inv}, \code{is-inv-has-inv},
  \code{id-is-inv}, \code{functor-pres-inv}\\
Lemma 35 & \code{FullyFaithful.ff-reflects-inv}\\
idtoiso is an equivalence & \code{FullyFaithful.rezk-idtoiso}\\
Lemma 37 & \code{FullyFaithful.ff-emb}\\
Remark 39 and converse & \code{FullyFaithful.ff-surj-iso}, \code{iso-ff}, \code{iso-surj}\\
\bottomrule
\end{tabular}
\end{center}
Composition is written in diagrammatic order in the files:
\code{hcomp g f} is $f\circ g$.

\subsection{Correspondence and differences}
\begin{itemize}
\item Each file begins with the informal proof given here, in more
  detail. A construction is locked (made opaque) right after the lemmas
  that characterise it. For example, \code{composite} is locked after
  properties (1) and (2) of Section~\ref{sec:comp}, and every later
  proof uses only those two properties.
\item The Agda proof postulates three facts: functors preserve
  composition, identities are invertible, and invertibility is a
  proposition. Here they are theorems (Proposition~\ref{prop:funct},
  Proposition~\ref{prop:inv}, and the definition). The three
  postulates of the Agda proof of Lemma 26 that were candidates for
  proof are proved as well: $\Hom_\II(0,1)$ is contractible, and the
  spine edges of $\Delta^2$ are its geometric faces.
\item Axiom~\ref{ax:rezk} is the parameter \code{rezk}. The equivalence
  with W\"arn's formulation (Remark~\ref{rem:rezk}) is proved here on
  paper but not formalised. Its ingredients are in the files
  (\code{mor-ends-equiv}, \code{is-inv-prop}, \code{fibrewise-from-total}).
\item The agreement of composition with the Segal witness is formalised
  for the underlying morphisms only, not for their endpoint proofs.
\item Surjectivity is propositionally truncated, using the truncation
  of the foundations.
\end{itemize}

\subsection{Figures}
Checking the whole directory takes 88 modules and 1674 declarations,
with no failures, in about 7\,s. Each file of
Sections~\ref{sec:segal}--\ref{sec:ff} checks in under 0.2\,s.

\subsection{Open}
Natural transformations, whiskering and the exchange law (W\"arn,
28--31), the Yoneda lemma, and full subcategories (Axiom 40 onwards)
are not yet treated. It is open whether any of them needs a coherence
of level 3.

\end{document}
